% Generated draft: requires mathematical and source review.
\documentclass[11pt]{article}
\usepackage[margin=25mm]{geometry}
\usepackage{fontspec}
\setmainfont{Latin Modern Roman}
\usepackage{amsmath,amssymb,mathtools,unicode-math}
\setmathfont{Latin Modern Math}
\usepackage{enumitem,needspace,xurl,hyperref}
\setlength{\parindent}{0pt}
\setlength{\parskip}{5pt}
\setlength{\emergencystretch}{3em}
\newcommand{\N}{\mathbb{N}}
\newcommand{\Z}{\mathbb{Z}}
\newcommand{\Q}{\mathbb{Q}}
\newcommand{\R}{\mathbb{R}}
\newcommand{\C}{\mathbb{C}}
\newcommand{\F}{\mathbb{F}}
\newcommand{\A}{\mathbb{A}}
\newcommand{\PP}{\mathbb{P}}
\newcommand{\E}{\mathbb{E}}
\newcommand{\eps}{\varepsilon}
\newcommand{\dd}{\,\mathrm{d}}
\begin{document}
\section*{UnsolvedMath ID: 700680}
\textbf{Source problem code:} catalogue entry
\subsection{Definitions and mathematical statement}
Let \(\Omega\) be an infinite set and let \(G\leq\operatorname{Sym}(\Omega)\) act \(2\)-transitively: for any \(\alpha,\beta,\gamma,\delta\in\Omega\) with \(\alpha\ne\beta\) and \(\gamma\ne\delta\), there exists \(g\in G\) such that \(g(\alpha)=\gamma\) and \(g(\beta)=\delta\). For distinct \(\alpha,\beta\in\Omega\), define the pointwise stabilizers
\[
G_\alpha=\{g\in G:g(\alpha)=\alpha\},\qquad
G_{\alpha\beta}=G_\alpha\cap G_\beta.
\]
Suppose there exists a subgroup \(A\leq G_\alpha\) such that
\[
G_\alpha=A\cdot G_{\alpha\beta}
\quad\text{and}\quad
\exists t\in G_{\alpha\beta}\ \forall a\in A,\quad tat^{-1}=a^{-1}.
\]
Here the displayed product is read as the set product \(\{ah:a\in A,\ h\in G_{\alpha\beta}\}\). The source does not explicitly state whether the notation additionally requires \(A\) to be normal in \(G_\alpha\), or the product to be a semidirect product; this qualification remains unresolved. No order condition on \(t\) is explicitly supplied.

\textbf{Classification problem.} Determine the structure of all infinite \(2\)-transitive permutation groups satisfying the preceding hypotheses, subject to clarification of the product notation.

\textbf{Additional question.} The source further supposes that \(A\setminus\{1\}\) ``contains at most two conjugacy classes of \(G_\alpha\).'' One precise possible interpretation is that the elements of \(A\setminus\{1\}\) represent at most two \(G_\alpha\)-conjugacy classes. Writing
\[
a\sim b\quad\Longleftrightarrow\quad
\exists h\in G_\alpha,\quad hah^{-1}=b
\qquad(a,b\in A\setminus\{1\}),
\]
this interpretation gives the condition
\[
\bigl|(A\setminus\{1\})/{\sim}\bigr|\leq 2.
\]
Under the intended conjugacy-class condition, must there exist a field \(K\) and a normal subgroup \(N\trianglelefteq G\) such that
\[
N\cong\operatorname{PSL}_2(K),\qquad
\operatorname{PSL}_2(K)=\operatorname{SL}_2(K)/Z(\operatorname{SL}_2(K))?
\]
Here \(\operatorname{SL}_2(K)\) is the group of determinant-one \(2\times2\) matrices over \(K\), and \(Z(H)\) denotes the center of a group \(H\). The source imposes no explicit characteristic restriction on \(K\).

The conjugacy-class wording needs review: it may instead require \(A\setminus\{1\}\) to be a union of at most two entire \(G_\alpha\)-conjugacy classes. This is stronger than the displayed interpretation unless the requisite conjugation invariance of \(A\) is established or assumed. A literal bound only on classes wholly contained in \(A\setminus\{1\}\) would be another, weaker reading. The source does not resolve these readings explicitly.
\subsection{Short English statement}
Classify the infinite \(2\)-transitive permutation groups whose point stabilizer is a product of a subgroup \(A\) and a two-point stabilizer containing an element that inverts every element of \(A\). Under the additional condition involving at most two point-stabilizer conjugacy classes in \(A\setminus\{1\}\), must the group contain a normal subgroup isomorphic to \(\operatorname{PSL}_2(K)\) for some field \(K\)? The meanings of the product notation and the conjugacy-class condition require clarification.
\subsection{Sources}
\begin{itemize}
\item Supplied local UnsolvedMath registry record: permanent ID \(700680\), exact source problem code \texttt{catalogue entry}, title \emph{Kourovka Notebook Problem 10.65}.
\item A. N. Fomin, Problem 10.65, \emph{The Kourovka Notebook (Unsolved Problems in Group Theory)}, 10th issue (1986), as attributed in the supplied record. The record also identifies the 21st issue edition (2026), \texttt{arXiv:1401.0300v48}, and names E. I. Khukhro and V. D. Mazurov as editors of the notebook collection.
\item The supplied registry labels the problem \emph{open} and records an update at \texttt{2026-10-16T14:00:00Z}. This is supplied registry information, not an independently verified assessment of the mathematical literature. The cited notebook text was not independently consulted for this draft.
\end{itemize}
\end{document}
